% back/notation.tex —— 符号表（后置材料）
% 依第 1 章 §1.1「基本记号」所列条目，并补入各章陆续引入的主要记号。
% 每条均在正文中有对应出现，不求穷尽，只收全书反复使用者；章次注于含义之后。
% 排版体例：与「术语索引」一致（longtable + booktabs，表头 \headfont）。
\cleardoublepage
\phantomsection
\chapter*{符号表}
\addcontentsline{toc}{chapter}{符号表}
\markboth{符号表}{符号表}

本书的记号在第 1 章开篇集中列出, 其余随各章内容陆续引入. 为便于查阅, 现将全书主要记号分类汇总于下; 含义之后括注该记号主要出现的章次, 未注者通贯全书.

\begingroup
\renewcommand{\arraystretch}{1.18}%
\begin{longtable}{@{}>{\raggedright\arraybackslash}p{0.27\textwidth}>{\raggedright\arraybackslash}p{0.63\textwidth}@{}}
  \toprule
  {\headfont 记号} & {\headfont 含义} \\
  \midrule
  \endfirsthead
  \toprule
  {\headfont 记号} & {\headfont 含义} \\
  \midrule
  \endhead
  \bottomrule
  \endlastfoot

  \multicolumn{2}{@{}l}{\headfont 一、数集与空间} \\*[3pt]
  $\R,\ \C$ & 实数域与复数域 \\
  $\R^n,\ \C^n$ & 实, 复 $n$ 维列向量空间 \\
  $\R^{m\times n},\ \C^{m\times n}$ & $m\times n$ 实, 复矩阵空间 \\
  $\ii$ & 虚数单位 $\sqrt{-1}$ \\
  $\mathcal{F}$ & 计算机浮点数集 (第 3 章) \\

  \addlinespace[6pt]
  \multicolumn{2}{@{}l}{\headfont 二、向量, 矩阵与基本运算} \\*[3pt]
  $A=[a_{ij}]$ & 第 $(i,j)$ 个元素为 $a_{ij}$ 的矩阵 \\
  $\mathbf{x},\mathbf{y},\mathbf{z}$ & 列向量, 以小写粗正体表示 \\
  $A^T,\ A^*$ & 转置, 共轭转置 \\
  $A^{-1},\ A^{\dagger}$ & 逆矩阵, Moore--Penrose 广义逆 (第 4 章) \\
  $I_n$ & $n$ 阶单位矩阵 \\
  $\mathbf{e}_i$ & 第 $i$ 个单位向量, 即 $I_n$ 的第 $i$ 列 \\
  $\mathbf{0}$ & 零矩阵或零向量, 尺寸视上下文而定 \\
  $|A|$ & 元素取绝对值的矩阵, $(|A|)_{ij}=|a_{ij}|$ \\
  $\det(A),\ \rank(A),\ \tr(A)$ & 行列式, 秩, 迹 \\
  $\dim(S)$ & 向量空间 $S$ 的维数 \\
  $\diag(\cdot)$ & 以给定元素为对角元的对角矩阵 \\
  $\spn\{\cdot\}$ & 由给定向量张成的子空间 \\
  $\operatorname{nullity}(A)$ & $A$ 的零度, 即 $\mathcal{N}(A)$ 的维数 (第 4 章) \\
  $A(i,j),\ A(i,:),\ A(:,j)$ & 仿 MATLAB 的元素, 行, 列记号 \\
  $A(k:l,p:q)$ & 由第 $k$ 至 $l$ 行, 第 $p$ 至 $q$ 列构成的子矩阵 \\

  \addlinespace[6pt]
  \multicolumn{2}{@{}l}{\headfont 三、范数与条件数} \\*[3pt]
  $\norm{\cdot}$ & 矩阵或向量的范数, 所指依上下文判定 \\
  $\norm{\cdot}_p$ & $p$ 范数 \\
  $\norm{\cdot}_1,\ \norm{\cdot}_2,\ \norm{\cdot}_\infty$ & 1 范数, 2 范数与 $\infty$ 范数 \\
  $\norm{A}_F$ & Frobenius 范数 \\
  $\triplenorm{\cdot}$ & 习题中表示矩阵范数的三竖线记号 \\
  $\kappa(A)$ & $A$ 的条件数, $\kappa(A)=\norm{A}\norm{A^{-1}}$ (第 3 章) \\
  $c(x)$ & 函数求值的条件数 (第 1 章) \\

  \addlinespace[6pt]
  \multicolumn{2}{@{}l}{\headfont 四、特征值, 奇异值与谱} \\*[3pt]
  $\lambda_i$ & 矩阵的第 $i$ 个特征值 \\
  $\Lambda$ & 以特征值为对角元的矩阵, $\Lambda=\diag(\lambda_1,\ldots,\lambda_n)$ \\
  $\rho(A)$ & $A$ 的谱半径, $\rho(A)=\max\{|\lambda|:\lambda\ \text{属于}\ A\ \text{的谱}\}$ (第 3 章) \\
  $\sigma_i$ & 第 $i$ 个奇异值 (第 8 章) \\
  $\Sigma,\ \Sigma_r$ & 奇异值构成的对角矩阵, $r=\rank(A)$ (第 8 章) \\
  $\mathbf{u}_i,\ \mathbf{v}_i$ & 第 $i$ 个左, 右奇异向量 (第 8 章) \\

  \addlinespace[6pt]
  \multicolumn{2}{@{}l}{\headfont 五、矩阵分解与标准形} \\*[3pt]
  $A=LU$ & LU 分解, $L$ 为下三角矩阵, $U$ 为上三角矩阵 (第 2 章) \\
  $A=QR$ & QR 分解, $Q$ 正交或酉, $R$ 为上三角矩阵 (第 4 章) \\
  $A=U\Sigma V^*$ & 奇异值分解 (第 4, 8 章) \\
  $Q^TAQ$ & 实对称矩阵的正交对角化, 即 $Q^TAQ=\diag(\lambda_1,\ldots,\lambda_n)$ (第 8 章) \\
  $A=X\Lambda X^{-1}$ & 可对角化矩阵的特征分解 (第 7 章) \\

  \addlinespace[6pt]
  \multicolumn{2}{@{}l}{\headfont 六、子空间与集合} \\*[3pt]
  $\mathcal{R}(A)$ & $A$ 的值域, 即列空间 (第 4 章) \\
  $\mathcal{N}(A)$ & $A$ 的零空间 (第 4 章) \\
  $\mathcal{S}^{\perp}$ & 子空间 $\mathcal{S}$ 的正交补 (第 4 章) \\
  $\mathcal{X}$ & 最小二乘问题的解集 (第 4 章) \\
  $\mathcal{P}_k$ & 次数不超过 $k$ 且 $P_k(0)=1$ 的多项式集合 (第 6 章) \\
  $\mathcal{M}_U$ & 可被酉矩阵 $U$ 同时对角化的矩阵集合 (第 6 章) \\
  $\mathcal{O}(\cdot)$ & 运算量或误差的阶 (大 $\mathcal{O}$ 记号) \\

  \addlinespace[6pt]
  \multicolumn{2}{@{}l}{\headfont 七、迭代法与 Krylov 子空间} \\*[3pt]
  $\mathbf{x}_k$ & 第 $k$ 步迭代向量 (第 5, 6 章) \\
  $\mathbf{r}_k$ & 第 $k$ 步残差, $\mathbf{r}_k=\mathbf{b}-A\mathbf{x}_k$ (第 6 章) \\
  $\mathcal{K}(A,\mathbf{r}_0,k)$ & 由 $A$ 与初始残差 $\mathbf{r}_0$ 生成的 $k$ 维 Krylov 子空间 (第 6 章) \\
  $\mathbf{r}_t,\ \mathbf{r}_p$ & 原方程组与预处理后方程组的当前残差 (第 10 章) \\
  $\rho$ & 重启 GMRES 方法的重启参数 (第 10 章) \\

  \addlinespace[6pt]
  \multicolumn{2}{@{}l}{\headfont 八、误差与浮点运算} \\*[3pt]
  $fl(x)$ & 实数 $x$ 的浮点数 (第 3 章) \\
  $\mu_\varepsilon$ & 机器精度 (第 3 章) \\
  $\epsilon$ & 小参数: 扰动幅度或残差相对精度 \\
  $\hat{\mathbf{x}}$ & 计算解, 即精确解的近似 (第 3 章) \\

  \addlinespace[6pt]
  \multicolumn{2}{@{}l}{\headfont 九、常微分方程与分数阶扩散方程} \\*[3pt]
  BVM,\ IVM & 边值法 (boundary value method) 与初值法 (initial value method) (第 9 章) \\
  DDE,\ LMF & 时滞微分方程 (delay differential equation) 与线性多步公式 (linear multistep formula) (第 9 章) \\
  $\mathcal{D}_{\nu_1,\nu_2}$ & BVM 的 $A_{\nu_1,\nu_2}$-稳定区域 (第 9 章) \\
  $\Gamma$ & 复平面中的一条曲线 (第 9 章) \\
  $\Gamma(z)$ & Gamma 函数, $\Gamma(z)=\int_0^\infty e^{-t}t^{z-1}\,\mathrm{d}t$ (第 10 章) \\
  ${}_{a}D_x^{\alpha},\ {}_{x}D_b^{\alpha}$ & $x$ 方向左, 右 Riemann--Liouville 分数阶导数 (第 10 章) \\
  ${}_{c}D_y^{\beta},\ {}_{y}D_d^{\beta}$ & $y$ 方向左, 右 Riemann--Liouville 分数阶导数 (第 10 章) \\
  $\alpha,\ \beta$ & 分数阶导数的阶, $1<\alpha,\beta<2$ (第 10 章) \\
  $\Omega,\ \partial\Omega$ & 求解区域及其边界 (第 10 章) \\
  $\tau$ & 时间步长, $\tau=T/M$ (第 10 章) \\
  $h_x,\ h_y$ & $x$, $y$ 方向的网格尺寸 (第 10 章) \\
  $N_1,\ N_2,\ M$ & 空间剖分数与时间剖分数 (第 10 章) \\
  $\underline{\mathbf{A}}^{(m+1)},\ \mathbf{u}^{(m+1)}$ & 第 $m+1$ 个时间层的系数矩阵与解向量 (第 10 章) \\

\end{longtable}
\endgroup
